% ============================================================================
% Fichier  : front/02_mode_emploi.tex
% Titre    : Mode d'emploi du cours
% Auteur   : MINKA MI NGUIDJOÏ Thierry Emmanuel
% Version  : 1.0 -- novembre 2026
% ============================================================================

\chapter*{Mode d'emploi du cours}
\addcontentsline{toc}{chapter}{Mode d'emploi du cours}
\markboth{Mode d'emploi}{Mode d'emploi}

\section*{Structure générale}

Le cours compte un prologue et six parties, soit 26~chapitres, pour
60~heures en présentiel réparties en trente séances de deux heures. Les
fiches d'auto-évaluation et les dix laboratoires de NDAP SUKULU sont des
travaux à faire hors de ces 60~heures.

\begin{tableaucours}{L{1.6cm}L{4.4cm}L{4.4cm}L{1.8cm}}{%
    \tableauHeader{Partie} &
    \tableauHeader{Thème} &
    \tableauHeader{Compétences visées} &
    \tableauHeader{Niveau}
}
Prologue & La lettre d'autorisation & Lire un mandat, séparer constatation
et opinion & \niveaubase \\
I & Le cadre du test & Métiers, droit, règles d'engagement, méthodes &
\niveaubase \\
II & Reconnaître & Sources ouvertes, balayage, vulnérabilités &
\niveaumoyen \\
III & Systèmes et réseaux & Exploitation, réseau, Active Directory,
privilèges, mots de passe & \niveaumoyen{} à \niveauexpert \\
IV & Applications et nuage & Web, injections, API, mobile, nuage &
\niveaumoyen{} à \niveauexpert \\
V & Au-delà des systèmes & Sans fil, ingénierie sociale, équipe rouge,
détection & \niveauexpert \\
VI & Restituer et progresser & Rapport, contre-test, divulgation,
laboratoire, mesure & \niveaumoyen \\
\end{tableaucours}

\section*{Les trente séances}

\begin{tableaucours}{L{1.4cm}L{8.4cm}L{2.4cm}}{%
    \tableauHeader{Séance} &
    \tableauHeader{Contenu} &
    \tableauHeader{Partie}
}
1 & Prologue~: la lettre d'autorisation & Prologue \\
2 & Le hacking éthique (Ch.~\ref{chap:ethique}) & I \\
3 & Le cadre juridique (Ch.~\ref{chap:droit}) & I \\
4 & Règles d'engagement (Ch.~\ref{chap:engagement}) & I \\
5 & Méthodologies et référentiels (Ch.~\ref{chap:methodes}) & I \\
6 & Reconnaissance passive (Ch.~\ref{chap:osint}) & II \\
7 & Balayage et énumération (Ch.~\ref{chap:balayage}) & II \\
8 & Analyser et prioriser (Ch.~\ref{chap:vulnerabilites}) & II \\
9 & L'exploitation (Ch.~\ref{chap:exploitation}) & III \\
10 & Attaques réseau (Ch.~\ref{chap:reseau}) & III \\
11--12 & Active Directory (Ch.~\ref{chap:ad}) & III \\
13 & Élévation de privilèges (Ch.~\ref{chap:privileges}) & III \\
14 & Post-exploitation (Ch.~\ref{chap:post}) & III \\
15 & Mots de passe (Ch.~\ref{chap:mdp}) & III \\
16--17 & Tester une application web (Ch.~\ref{chap:web}) & IV \\
18 & Injections et contrôle d'accès (Ch.~\ref{chap:injections}) & IV \\
19 & API et mobile (Ch.~\ref{chap:api}) & IV \\
20 & Nuage et conteneurs (Ch.~\ref{chap:nuage}) & IV \\
21 & Sans fil et objets connectés (Ch.~\ref{chap:sansfil}) & V \\
22 & Ingénierie sociale (Ch.~\ref{chap:social}) & V \\
23--24 & Équipe rouge et équipe violette (Ch.~\ref{chap:rouge}) & V \\
25 & Détection et contournement (Ch.~\ref{chap:detection}) & V \\
26 & Le rapport de test (Ch.~\ref{chap:rapport}) & VI \\
27 & Contre-test (Ch.~\ref{chap:retest}) & VI \\
28 & Divulgation coordonnée (Ch.~\ref{chap:divulgation}) & VI \\
29 & Le laboratoire (Ch.~\ref{chap:labo}) & VI \\
30 & Mesurer~; synthèse (Ch.~\ref{chap:mesurer}) & VI \\
\end{tableaucours}

\section*{Les encadrés}

\begin{boxdef}
\textbf{Définition}. Introduit un terme avec sa définition~; les termes
définis sont indexés et repris dans le glossaire (Annexe~\ref{ann:glossaire}).
\end{boxdef}

\begin{boxscenario}{Titre du scénario}
\textbf{Scénario opérationnel}. Un épisode d'une des opérations du
cours~: LOGONE (test d'intrusion de LiZbiZ), MBAM (équipe rouge de la BANQUE
SANAGA), MEFOU (portail et API de l'ANPP).
\end{boxscenario}

\begin{boxobjet}{T-xx · Intitulé}
\textbf{Fiche de constat}. Présente un constat de test sous une forme
constante~: description, preuve, gravité, impact, recommandation,
détection.
\end{boxobjet}

\begin{boxprocedure}
\textbf{Procédure de test}. Dit, pas à pas, ce qu'il faut faire, dans le
cadre d'une autorisation écrite.
\end{boxprocedure}

\begin{boxmetier}{Pentesteur}
\textbf{Application métier}. Traduit le chapitre en gestes concrets pour
l'un des métiers visés~: pentesteur, auditeur de sécurité, analyste SOC et
équipe rouge, RSSI.
\end{boxmetier}

\begin{boxjuris}
Cite un texte de loi d'après sa version officielle, puis en tire la
portée pour le testeur.
\end{boxjuris}

\begin{boxalert}
\textbf{Violation typique} et \textbf{règle}~: une erreur fréquente, et ce
qu'il faut faire à la place.
\end{boxalert}

\begin{boxexo}
\textbf{Exercice}. Dix par chapitre, gradués en trois niveaux~:
\niveaubase\ fondamental, \niveaumoyen\ intermédiaire, \niveauexpert\
expert. Les corrigés figurent dans le livret de correction.
\end{boxexo}

\section*{Conventions}

\begin{itemize}
    \item Toute technique est présentée pour des systèmes que l'on est
          autorisé à tester, avec sa détection et sa parade.
    \item Chaque analyse indique un niveau de confiance~: faible, moyenne ou
          élevée, et reste séparée de la constatation.
    \item Les adresses IP et les noms de domaine des exemples sont fictifs ou
          réservés à la documentation.
    \item Les \termecle{termes clés} apparaissent en gras~; les normes
          s'écrivent \norme{ISO/IEC~29147}, les lois camerounaises
          \loicam{n°~2010/012}.
\end{itemize}

\section*{Matériel d'accompagnement}

\begin{itemize}
    \item \textbf{Annexes}~: glossaire, outils, modèles et procédures,
          textes juridiques et normatifs, bibliographie commentée, contacts
          (Annexes~\ref{ann:glossaire} à~\ref{ann:contacts}).
    \item \textbf{Livret de correction}~: corrigés des exercices et des
          fiches, réservé au formateur.
    \item \textbf{NDAP SUKULU}~: fiches d'auto-évaluation et laboratoires,
          présentés dans la section qui suit.
\end{itemize}
